\section{Existential tests for controlled families}
\label{sec:existential}

The polynomial operations developed above allow us to construct much
larger polynomials than we can write down.  We now show that real
feasibility remains decidable in a fixed counting level for these
families.  This will let us apply an existential test after the
coefficient and determinant constructions used for the two quantifier
blocks.

\begin{theorem}\label{thm:controlled-existential}
Suppose that an input word specifies a conjunction of polynomially many
equations and strict inequalities
\[
 A_i(x)=0,\qquad B_j(x)>0,
\]
in polynomially many real variables.  Assume that the numbers of
variables and tests, and the type of each addressed test, are computable
in polynomial time, and that the polynomials form controlled families
uniformly in their indices.  Then the existence of a real solution is
an available decision.  Its counting level depends only on the fixed
family constructions and their evaluation algorithms, not on the input
or any short outer indices.
\end{theorem}

We begin the proof by replacing each inequality $B_j(x)>0$ by
$B_j(x)z_j^2-1=0$, with a new real variable $z_j$.  This replacement is
equivalent over $\R$.  The sum of the squares of all equation left sides
is a controlled polynomial $F$ by
Proposition~\ref{prop:controlled-closure}.  It is nonnegative on
$\R^n$, and
\[
 F\text{ has a real zero}
 \quad\Longleftrightarrow\quad
 \text{the original conjunction is feasible}.
\]
Add an unused variable if necessary so that $n\ge1$; an empty conjunction
gives $F=0$.  It therefore suffices to prove the theorem for the zero set
of a controlled nonnegative polynomial.

We first turn this zero-set question into a question about one
univariate polynomial.  A coercive perturbation supplies a minimum
whose boundedness detects a zero of $F$: compact sublevel sets will
supply a convergent subsequence whose limit is a zero.  Its critical-value
polynomial then replaces the limiting question by one large integer
specialization.  Finally, the root labels from
Proposition~\ref{prop:root-labels} convert the univariate question into
the sign of an integer determinant.

\subsection{A minimum that detects a zero}

Choose an odd integer $\ell\ge3$ larger than the degree bound for $F$.
It has a polynomial-length binary encoding and is computable in
polynomial time.  For $1\le i\le n$, put $a_i=(4\ell)^i$ and define
\begin{equation}\label{eq:existential-deformation}
 \begin{split}
 G(x)&=\sum_{i=1}^n
   \bigl(x_i^{\ell+1}-(\ell+1)a_i^\ell x_i
                    +(\ell+1)a_i^{\ell+1}\bigr),\\
 h_u(x)&=G(x)+(\ell+1)uF(x),\\
 r_i(x,u)&=a_i^\ell-u\,\partial_iF(x),\qquad
 f_i(x,u)=x_i^\ell-r_i(x,u).
 \end{split}
\end{equation}
Coercive perturbations and critical-point equations are classical tools
in real algebraic geometry; see, for example,
\citet{BasuPollackRoy2006}.  Here the separated constants $a_i$ serve an
additional purpose: the critical values at $u=0$ will all be distinct.

All polynomials in Equation~\eqref{eq:existential-deformation}, including
$u$ as an indeterminate, are controlled.  Indeed, the $a_i$ themselves
have polynomial bit length, and their displayed powers are allowed by
Proposition~\ref{prop:controlled-closure}.  Moreover,
$\deg_x r_i<\ell$.  Apply
Proposition~\ref{prop:multiplication-determinant} to the relations
$f_1,\ldots,f_n$ and to multiplication by $h_u$, obtaining
\begin{equation}\label{eq:existential-characteristic}
 Q(u,y)=\det(yI-M_{h_u})\in\Z[u,y],
 \qquad D=\ell^n.
\end{equation}
This is a controlled polynomial, monic of degree $D$ in $y$.

Each summand of $G$ has derivative
$(\ell+1)(x_i^\ell-a_i^\ell)$.  Since $\ell$ is odd, its unique real
minimum occurs at $x_i=a_i$ and equals $a_i^{\ell+1}$.  The even leading
power makes $G$ coercive: $G(x)\to+\infty$ as
$\lVert x\rVert\to\infty$.  In particular,
\begin{equation}\label{eq:existential-coercivity}
 h_u(x)\ge G(x)\ge\sum_{i=1}^n a_i^{\ell+1}>0
 \qquad(u\ge0).
\end{equation}

\begin{lemma}\label{lem:bounded-minimum}
For the nonnegative polynomial $F$ and deformation in
Equation~\eqref{eq:existential-deformation}, the minimum
\[
 m(u)=\min_{x\in\R^n}h_u(x)
\]
exists for every real $u\ge0$.  The function $m$ is continuous and
nondecreasing, and $Q(u,m(u))=0$.  Furthermore,
\[
 F\text{ has a real zero}
 \quad\Longleftrightarrow\quad
 m(u)\text{ remains bounded as }u\to+\infty.
\]
\end{lemma}

\begin{proof}
Existence follows from continuity and coercivity.  Since $F\ge0$, each
$h_u(x)$ is nondecreasing in $u$, and so is its minimum.  At a minimizer,
\[
 \partial_i h_u=(\ell+1)f_i=0\qquad(1\le i\le n).
\]
Evaluation at this common zero of the relations shows that the value of
$h_u$ is a root of its characteristic polynomial, by
Lemma~\ref{lem:real-eigenline}.

For continuity, fix a compact interval $I\subseteq[0,\infty)$.  Every
minimizer for a parameter $u\in I$ lies in the sublevel set
\[
 G(x)\le h_u(x)=m(u)\le h_u(0)
 \le\max_{v\in I}h_v(0).
\]
This set is compact and contains all these minimizers.  On it the
functions $h_u$ vary uniformly continuously with $u$.  The difference
of their minima is at most their uniform difference there, proving
continuity of $m$ on $I$.

If $F(x_0)=0$, then $m(u)\le G(x_0)$ for every $u\ge0$.  Conversely,
suppose that $m(u)\le M$ for all $u\ge0$, and take $u_k\to+\infty$
with minimizers $x_k$.  Equation~\eqref{eq:existential-coercivity} gives
\[
 G(x_k)\le M,
 \qquad 0\le F(x_k)\le\frac{M}{(\ell+1)u_k}.
\]
A subsequence converges in the compact sublevel set of $G$ to a point
$x_*$, and continuity yields $F(x_*)=0$.
\end{proof}

\subsection{Replacing the limit by one specialization}

The boundedness criterion is not yet a finite test.  We will choose a
value cutoff $R$ and a parameter $U$ so that a bounded minimum lies below
$R$, whereas a divergent minimum exceeds $R$ by parameter $U$.  A
finite limit of $m(u)$ must be a root of the leading coefficient of $Q$
as a polynomial in $u$, itself a polynomial in $y$, so a coefficient
bound supplies $R$.  Once $R$ is chosen, taking $U$ larger
than every positive real root of the nonzero polynomial $Q(u,R)$
prevents the minimum from crossing $R$ after $U$.  We also take $U$
beyond every parameter where $Q(u,y)$ has a repeated root:
Lemma~\ref{lem:real-eigenline} will then realize every real root of
$Q(U,y)$ as a value at a real critical point.  We first bound those
exceptional parameters.

\begin{lemma}\label{lem:existential-discriminant}
The polynomial $Q(0,y)$ has $D$ distinct complex roots.  If
$\norm{Q}\le2^H$, where $H\ge0$ is an integer, there is a nonzero
polynomial $\Delta\in\Z[u]$ such that every parameter where $Q(u,y)$
has a repeated complex root is a zero of $\Delta$, and
\begin{equation}\label{eq:existential-discriminant-bound}
 \norm{\Delta}\le2^{H_\Delta},\qquad
 H_\Delta=(2D)^2+2D(H+D+1).
\end{equation}
\end{lemma}

\begin{proof}
At $u=0$, the common zeros of the relations are
$x_i=a_i\zeta_i$, where $\zeta_i^\ell=1$.  Their values under $G$ are
\begin{equation}\label{eq:existential-unperturbed-values}
 \sum_{i=1}^n a_i^{\ell+1}
             \bigl((\ell+1)-\ell\zeta_i\bigr).
\end{equation}
They are pairwise distinct.  To see this, let $j$ be the largest index
where two root-of-unity tuples differ, and put
$W_i=a_i^{\ell+1}$.  The distance between distinct $\ell$th roots of
unity is at least
$2\sin(\pi/\ell)\ge4/\ell$.  Thus the contribution at index $j$ to
the difference in Equation~\eqref{eq:existential-unperturbed-values}
has modulus at least $4W_j$.  The sum of the earlier contributions has
modulus at most
\[
 2\ell\sum_{i<j}W_i
 <\frac{2\ell}{(4\ell)^{\ell+1}-1}W_j
 <4W_j.
\]
The contribution at index $j$ cannot cancel.  By
Lemma~\ref{lem:real-eigenline}, all $\ell^n=D$ distinct values are roots
of the degree-$D$ polynomial $Q(0,y)$.

Let $\Delta(u)$ be the determinant, in the monomial bases, of the square
coefficient map
\begin{equation}\label{eq:existential-sylvester-map}
 (B_1,B_2)\longmapsto B_1Q+B_2Q_y,
 \quad \deg_y B_1<D-1,\quad \deg_y B_2<D,
\end{equation}
whose target consists of polynomials of degree less than $2D-1$.
At any repeated root of $Q$, evaluation annihilates the image, so the
specialized determinant vanishes.  At $u=0$ the map is injective:
evaluating a zero image at the $D$ distinct roots of $Q(0,y)$ gives
$B_2=0$, because $Q_y$ is nonzero at those roots and $\deg_y B_2<D$;
then $B_1=0$.  Consequently $\Delta(0)\ne0$.

The map has size $2D-1$, and each entry, as a polynomial in $u$, has
coefficient norm at most $D2^H$.  Determinant expansion gives
\[
 \norm{\Delta}\le(2D-1)!(D2^H)^{2D-1}
 \le2^{(2D)^2+2D(H+D+1)},
\]
as asserted.
\end{proof}

Recall the elementary root bound: a nonzero integer polynomial of
coefficient norm at most $2^b$ has all its complex roots in
$|z|\le1+2^b$.  For $|z|>1+2^b$, its leading coefficient has modulus
at least $1$ and the geometric-series estimate makes its leading term
strictly larger than the sum of the lower terms.

Choose a polynomial-time computable short bound $H$ with
$\norm{Q}\le2^H$, as supplied by
Proposition~\ref{prop:multiplication-determinant}, and let $V$ be a
short bound on $\deg_uQ$.  Define
\begin{equation}\label{eq:existential-specialization-parameters}
 \begin{aligned}
 \rho&=H+4,& R&=2^\rho,\\
 \sigma&=H_\Delta+H+D\rho+4,& U&=2^\sigma,\\
 K&=H+\sigma V.
 \end{aligned}
\end{equation}
The exponents and bounds in this display are short integers.  The
integers $R$ and $U$ are used as controlled constant families; their
binary expansions need not have polynomial length.

\begin{proposition}\label{prop:specialization}
Let $F\in\Z[x_1,\ldots,x_n]$ be nonnegative on $\R^n$, let
$n\ge1$, and choose an odd $\ell\ge3$ with $\ell>\deg F$.
Define $G,h_u,f_i,Q,D$ by
Equations~\eqref{eq:existential-deformation} and
\eqref{eq:existential-characteristic}.  Let $H,V$ be nonnegative integers bounding
$\log_2\norm{Q}$ and $\deg_uQ$, respectively, and choose the parameters
in Equation~\eqref{eq:existential-specialization-parameters}.  Then
\[
 q(y)=Q(U,y)
\]
is monic and squarefree, has degree $D$ and coefficient norm at most
$2^K$, and satisfies
\begin{equation}\label{eq:existential-root-equivalence}
 F\text{ has a real zero}
 \quad\Longleftrightarrow\quad
 q\text{ has a real root less than }R.
\end{equation}
Every real root of $q$ is the value of $h_U$ at a real common zero of
$f_1(\cdot,U),\ldots,f_n(\cdot,U)$.

When $F$ varies in a controlled family, $\ell$ is chosen from its
degree bound as above, and $H,V$ are polynomial-time computable short
bounds, $q$ is controlled.  All remaining bound parameters in
Equation~\eqref{eq:existential-specialization-parameters}, including
$D,H_\Delta$, are computable in polynomial time and have polynomial
binary length.
\end{proposition}

\begin{proof}
Write
\[
 Q(u,y)=\sum_{j=0}^{a} A_j(y)u^j,\qquad A_a\ne0.
\]
The polynomial $A_a$ has integer coefficients and norm at most $2^H$.
The root bound and the choice of $R$ give $A_a(R)\ne0$.  Hence
$Q(u,R)$ is a nonzero integer polynomial in $u$, with norm at most
$2^H R^D=2^{H+D\rho}$.  Applying the root bound to this polynomial and
to $\Delta$ from Lemma~\ref{lem:existential-discriminant} yields
\begin{equation}\label{eq:existential-no-crossing}
 Q(u,R)\ne0,\qquad \Delta(u)\ne0
 \qquad\text{for every real }u\ge U.
\end{equation}
In particular, $q$ is squarefree.  Monicity and degree are preserved by
specialization, and $\norm{q}\le2^H U^V=2^K$.  Each real root is simple,
so Lemma~\ref{lem:real-eigenline} realizes it at a real common zero of
the specialized relations.

If $F$ has a real zero, Lemma~\ref{lem:bounded-minimum} gives a finite
limit $\mu=\lim_{u\to\infty}m(u)$.  Divide $Q(u,m(u))=0$ by $u^a$
and take this limit.  Every lower-power term tends to zero, since
$m(u)$ is bounded, and therefore $A_a(\mu)=0$.  The root bound gives
\[
 m(U)\le\mu\le|\mu|\le1+2^H<R.
\]
Thus $m(U)$ is a real root of $q$ below $R$.

If $F$ has no real zero, Lemma~\ref{lem:bounded-minimum} and monotonicity
give $m(u)\to+\infty$.  Its continuity and
Equation~\eqref{eq:existential-no-crossing} imply $m(U)>R$: otherwise
the minimum would equal $R$ at some parameter $u\ge U$.
Every real root of $q$ is a value of $h_U$ at a real point, and is
therefore at least $m(U)>R$.  This proves
Equation~\eqref{eq:existential-root-equivalence}.

Finally, $D=\ell^n$ has polynomial binary length because $n$ is
polynomially bounded and $\ell$ is short.  The determinant rule supplies
the short bounds $H,V$ in polynomial time.  The remaining exponents
are obtained by a fixed number of sums and products of short integers.
Substituting the controlled constant $U=2^\sigma$ into $Q$ preserves
control by Proposition~\ref{prop:controlled-closure}.  Neither the
highest nonzero index $a$ nor the coefficients of $\Delta$ need to be
computed: they justify the specialization, whose definition uses only
the stated bounds.
\end{proof}

\subsection{One root label and one determinant sign}

Proposition~\ref{prop:specialization} has reduced feasibility to the
existence of a real root of the controlled polynomial $q$ below $R$.
Apply Proposition~\ref{prop:root-labels} to $q$ with
$b_0(y)=R-y$.  It supplies a family of integer polynomials $L$, indexed
by short labels, with the following properties at real roots $\alpha$
of $q$:
\begin{itemize}
 \item $L(\alpha)\ge\tfrac12$ or $L(\alpha)\le-\tfrac12$;
 \item $L(\alpha)>0$ implies $\alpha<R$;
 \item each real root $\alpha<R$ has a label whose polynomial is
       positive at that root and negative at every other real root.
\end{itemize}
The family is controlled uniformly in the label.  Multiplying $L$ by
$-4D$ gives a sign margin large enough that a small integer shift
preserves all real-root signs.  We use that shift to avoid zero factors
at nonreal roots.  For a short choice $e\in\{0,\ldots,D\}$, define
\begin{equation}\label{eq:existential-sign-determinant}
 g(y)=-4DL(y)+e,
 \qquad X=\det M_{g(h_U)}.
\end{equation}
The multiplication matrices here use the relations specialized at
$u=U$.  They respect polynomial composition, and hence
\begin{equation}\label{eq:existential-product-identity}
 X=\det g(M_{h_U})
  =\prod_{q(\alpha)=0}g(\alpha)\in\Z.
\end{equation}
The product runs over all complex roots of $q$.

At a real root, the label margin gives
\[
 L(\alpha)>0\ \Longrightarrow\ g(\alpha)\le-D<0,
 \qquad
 L(\alpha)<0\ \Longrightarrow\ g(\alpha)\ge2D>0.
\]
Since $g$ has real coefficients, a conjugate nonreal pair contributes
$g(\alpha)g(\overline\alpha)=|g(\alpha)|^2\ge0$ to
Equation~\eqref{eq:existential-product-identity}.  Consequently $X<0$
implies that some real root has $g(\alpha)<0$, and such a root lies
below $R$.

Conversely, if $q$ has a real root below $R$, choose a label positive
at exactly that real root.  Since $q$ is squarefree, every
$e\in\{0,\ldots,D\}$ then gives
exactly one negative real factor and all other real factors positive.
At any fixed nonreal root, $g(\alpha)=0$ excludes at most one value of
$e$.  Since there are at most $D$ such roots and $D+1$ available
choices, some $e$ leaves every nonreal factor nonzero.  Each conjugate
pair now contributes a positive factor, so $X<0$.  We have proved
\begin{equation}\label{eq:existential-final-criterion}
 F\text{ has a real zero}
 \quad\Longleftrightarrow\quad
 X<0\text{ for some short label and some }0\le e\le D.
\end{equation}

\begin{proof}[Completion of the proof of
Theorem~\ref{thm:controlled-existential}]
The specialized relations and $g(h_U)$ are controlled by the closure
rules.  A second application of
Proposition~\ref{prop:multiplication-determinant} therefore makes $X$
a controlled integer family: it is $(-1)^D$ times the constant
coefficient of the characteristic polynomial of $M_{g(h_U)}$.
Corollary~\ref{cor:controlled-sign} gives its sign as an available
function, uniformly in the input, the label, and $e$.  Short-index
existential quantification in
Equation~\eqref{eq:existential-final-criterion} is available as well.

Every bound and prime cutoff is computed from the preceding controlled
families by the fixed rules of
Proposition~\ref{prop:controlled-closure} and the determinant rule.
The construction uses two applications of the determinant rule, one
root-label construction, and a fixed number of closure operations.
The possibly large degrees, exponents, and summation ranges all have
short binary encodings; none causes an input-dependent nesting of
counting computations.  Thus the resulting level is fixed uniformly
over the original inputs and any short outer indices, as required.
\end{proof}
